% !TeX root = ../main.tex
\chapter{Lorentzian causality}\label{f:causality}

The Levi--Civita connection preserves squared speed. For a positive-definite metric this controls length; for a Lorentzian metric it also distinguishes directions that can carry a signal from directions that cannot. The distinction starts in one tangent space. Whether a signal can return to its starting event, or whether there is a surface that every signal must cross, is a question about whole paths.

\section{Choosing a future cone}

The inertia of a real nondegenerate symmetric form is its number of negative and positive squares in a diagonal basis. These numbers are independent of that basis by Sylvester's law of inertia; for a smooth nondegenerate metric they are locally constant, since the eigenvalues of its symmetric matrix in a local frame vary continuously and none can cross zero. A Lorentzian metric in dimension $n\geq2$ has one negative and $n-1$ positive squares. At a point, choose a basis in which $g(v,v)=-(v^0)^2+|\mathbf v|^2$. A nonzero vector is timelike, null, or spacelike according as this quantity is negative, zero, or positive. A causal vector is timelike or null. Timelike vectors have two components, $v^0>|\mathbf v|$ and $v^0<-|\mathbf v|$. These are convex cones: addition preserves the strict inequalities by the triangle inequality. Each closed cone also includes its null boundary and zero.

A \emph{time orientation} continuously selects one of the timelike components as future; nonzero vectors in its closure are the future causal vectors. Such a choice is equivalent to a smooth timelike vector field $D$. Indeed a field selects the component containing $D_x$ at each $x$. Conversely choose local smooth timelike fields in the selected components and sum them with a partition of unity. At each point the terms with positive weights all lie in the same open convex cone, so the sum is timelike, not zero. In a basis with $D$ along the positive time axis, $g(D,v)<0$ for every nonzero future causal $v$ and has the opposite sign for a past causal $v$. This supplies a test that does not depend on the basis. Parallel transport from Chapter~\ref{f:connections} along a path in a time-oriented manifold preserves the future timelike cone: it preserves the norm, and the sign of the pairing with $D$ cannot change without a zero. By continuity it preserves the future null directions too. A causal path, however, need not have parallel velocity.

\begin{exercise}[A time direction that does not close]
The isometry $(t,x)\mapsto(-t,x+1)$ of the Minkowski plane generates a free, properly discontinuous action. Is its smooth quotient time-orientable? Compare with Example~\ref{f:circle-line}.

\emph{Solution.} A time orientation would pull back to an invariant choice on the connected covering plane. There it must choose $v^0>0$ everywhere or $v^0<0$ everywhere for its timelike directions; the generator exchanges the choices. The subgroup generated by the square acts by $(t,x)\mapsto(t,x+2)$ and preserves them. It is the choice of time cone that twists, just as the choice of real unit frame twisted on the circle.
\end{exercise}

\section{Reachability by directed paths}

A future timelike path may be taken piecewise $C^1$ with future timelike velocities on each regular piece. For causal paths we allow locally Lipschitz curves with nonzero future causal velocity almost everywhere, up to increasing reparametrization; constant waiting intervals can be suppressed. This class does not require a velocity at a corner. Write $p\ll q$ if a future timelike path runs from $p$ to $q$, and $p\leq q$ if $p=q$ or such a causal path exists. Set
\[
 I^+(p)=\{q:p\ll q\},\qquad J^+(p)=\{q:p\leq q\};
\]
these are the chronological and causal futures; reversing the direction gives $I^-(p),J^-(p)$. For a set, take the union of the corresponding futures of its points. Concatenation gives transitivity. A chronological spacetime has no closed timelike path; a causal spacetime has no nonconstant closed causal path, so its relation $\leq$ is a partial order.

In a convex normal neighborhood $U$ of $p$, write $v=\exp_p^{-1}(q)$. Future timelike $v$ gives a future timelike radial geodesic to $q$, and future null $v$ a null one. The local causal comparison theorem (the Lorentzian Gauss lemma and its consequence for curves in $U$, \cite[Ch.~5]{ON}) gives the converses: for $q\ne p$, $q\in I^+(p;U)$ exactly for future timelike $v$, and $q\in J^+(p;U)$ exactly for future causal $v$. Thus inside this neighborhood the sign of $\sigma_g(q,p)$ from \eqref{f:world-function} tests causal character, but not future versus past. The notation $;U$ matters: a path permitted to leave $U$ need not satisfy this endpoint test.

The chronological future is open. On the last smooth part of a timelike path ending at $q$, pick a point $r$ sufficiently near $q$ that this part stays in a convex normal neighborhood of $q$ containing $r$. Then $\exp_r^{-1}(q)$ is future timelike by local comparison. Initial velocities to nearby $q'$ remain in this open cone, giving a new final segment to append to the path ending at $r$. In the Minkowski plane $I^+(p)$ is given by $\Delta t>|\Delta x|$ and $J^+(p)$ by $\Delta t\geq|\Delta x|$: integrate the velocity inequalities for necessity and use a straight segment for sufficiency. Thus the causal future need not be open. Nor must it be closed. Remove $(0,0)$ from this plane and let $p=(-1,-1)$ and $q=(1,1)$, in $(t,x)$ order. Any future causal path in these coordinates satisfies $|\dot x|\leq\dot t$. Reaching $q$ would saturate the integrated bound, forcing $\dot x=\dot t$ almost everywhere and hence $x=t$ throughout the path. The missing origin forbids this. Yet $(1+\epsilon,1)$ for $\epsilon>0$ can be reached by a straight timelike segment avoiding the origin. In this spacetime a limit of points in $I^+(p)$ need not even belong to $J^+(p)$.

\section{Cauchy surfaces and global hyperbolicity}

An inextendible causal curve has neither a future nor a past extension in the spacetime, even after reparametrization. An infinite parameter interval is not the test: a curve approaching an ordinary point in finite coordinate time could be made to take infinite parameter time simply by relabelling it. A smooth spacelike hypersurface is \emph{Cauchy} if every inextendible causal curve meets it exactly once. A \emph{time function} is a continuous function increasing strictly along future causal paths. A \emph{temporal function} $T$ is smooth with everywhere past timelike gradient; $\mathrm dT(v)=g(\operatorname{grad}T,v)>0$ for future causal $v\ne0$, so it is a time function. A temporal function onto $\mathbb R$ is \emph{Cauchy temporal} when all its levels are Cauchy hypersurfaces; equivalently, its values along every inextendible causal curve exhaust $\mathbb R$.

Call a neighborhood causally convex if every causal segment with both endpoints in it stays in it. \emph{Strong causality} means there are arbitrarily small such neighborhoods at every point. Unlike geodesic convexity, this tests all directed paths in the spacetime, not just geodesics inside a chosen open set. A smooth time-oriented spacetime is \emph{globally hyperbolic} if it is strongly causal and every causal diamond $J^+(p)\cap J^-(q)$ is compact. Compactness prevents a sequence of possible intermediate events from escaping the spacetime while both endpoints stay fixed.

\begin{theorem}[Cauchy characterization and smooth splitting]\label{f:cauchy-splitting}
Let $M$ be a connected smooth time-oriented Lorentzian manifold without boundary. Global hyperbolicity, existence of a smooth spacelike Cauchy hypersurface, and existence of a Cauchy temporal function onto $\mathbb R$ are equivalent. In this case one can choose coordinates on a product $M\simeq\mathbb R\times\Sigma$ in which
\[
 g=-\beta(T,y)\,\mathrm dT^2+h_T(y),\qquad \beta>0 ,
\]
with $h_T$ a smooth family of Riemannian metrics and each slice a spacelike Cauchy hypersurface.
\end{theorem}

The equivalence and the existence of a \emph{smooth} Cauchy temporal function are global theorems, not consequences of diagonalizing the metric \cite{BGP,BS}. Once $T$ is available, the reason for the product formula is instructive. Flow along $D=\operatorname{grad}T/g(\operatorname{grad}T,\operatorname{grad}T)$; then $D(T)=1$. A maximal flow line cannot stop at a finite value of $T$: if it had an endpoint, the smooth ODE would extend it, and without an endpoint it would be future- or past-inextendible missing later or earlier Cauchy levels. Hence every line runs through $T=0$ and all other levels. The flow and its inverse (flow back by time $T$) give the product diffeomorphism. Its time vector is normal to the slices, which removes the mixed metric terms; the theorem does not say that $\beta=1$ or that $h_T$ is independent of $T$.

\section{Hypersurfaces: normals, screens and transversals}

Let a smooth hypersurface $F$ be given locally by $r=0$ with $\nu=\mathrm dr\ne0$. Write $k=\nu^\sharp$ for the vector characterized by $g(k,X)=\nu(X)$. Then $T_pF=k^\perp$, and the radical of the restricted metric, the subspace orthogonal to all of $T_pF$ inside $T_pF$, is $T_pF\cap\mathbb R k$. If $g^{-1}(\nu,\nu)<0$, the removed direction is timelike and $g|_{T_pF}$ is positive definite: the hypersurface is spacelike there. A positive conormal norm instead gives a timelike hypersurface, with one negative direction left in $T_pF$. In either case $k$ is transverse and may be normalized. The conclusions are unchanged on multiplying $r$ by a nowhere-zero smooth factor. At a corner this calculation classifies each incident smooth face separately.

\begin{proposition}[The null screen]\label{f:null-screen}
At a null point of $F$, $k$ is tangent, the radical is $\mathbb R k$, and $T_pF/\mathbb R k$ has a canonical positive-definite metric. A screen, meaning a spacelike complement to this line inside $T_pF$, realizes that quotient but need not be unique. A transverse vector lies outside $T_pF$, not in a screen.
\end{proposition}
\begin{proof}
Because $\nu(k)=0$, the radical computation gives the line $\mathbb R k$. Choose $w$ with $g(k,w)=-1$, which is possible by nondegeneracy. Put $l=w+\tfrac12 g(w,w)k$. Then $g(l,l)=0$ and $g(k,l)=-1$, so $l$ is transverse and the plane of $k,l$ has one positive and one negative square. Its orthogonal complement $S$ is positive definite. We have $T_pF=\mathbb R k\oplus S$; changing a representative by a multiple of $k$ does not alter its pairing with another tangent vector. Thus the quotient metric is independent of the choice of $l$.
\end{proof}

On a face that is null everywhere, the radical lines are tangent to null geodesics up to a change of parameter. Indeed $k=\operatorname{grad}r$, so symmetry of the Hessian gives $g(\nabla_k k,Y)=g(\nabla_Y k,k)=\tfrac12 Yg(k,k)$; the Hessian is symmetric because the Levi--Civita connection has zero torsion. For $Y$ tangent to $F$ the last derivative vanishes there. Hence $\nabla_k k$ is a multiple of $k$, and an integrating factor along each integral curve rescales its tangent to a parallel one. This does not turn $k$ into a transverse normal. For example, with $g=-\mathrm du\otimes\mathrm dv-\mathrm dv\otimes\mathrm du+\sum_A(\mathrm dy^A)^2$, on the wall $u=0$ of $u\geq0$ one has $k=-\partial_v$. The field $-\partial_u$ is an outward null transversal, while the $\partial_{y^A}$ span a screen. The normal line $TM|_F/TF$ from Chapter~\ref{f:submanifolds} is a different quotient from the screen quotient $TF/\mathbb R k$: the transversal has nonzero class in the former, whereas $k$ has zero class there. The outward-first orientation uses a transversal even when the metric normal is tangent.

\begin{exercise}[A face of changing type]
In the Minkowski plane classify the graph $t=x^2$ using its conormal. Is there an interval of null generators on it?

\emph{Solution.} The conormal $\mathrm dt-2x\,\mathrm dx$ has norm squared $4x^2-1$. The graph is spacelike for $|x|<1/2$, timelike for $|x|>1/2$, and null just at $x=\pm1/2$. The conclusion about null geodesic generators requires a null hypersurface on an open part of the face, not merely a null tangent at an isolated point.
\end{exercise}

\section{A cylinder with a time shift}

Consider the Minkowski plane with $g=-\mathrm dt^2+\mathrm dx^2$ and future $\partial_t$. Fix $L>0$ and divide by the translations generated by $(t,x)\mapsto(t+\tau,x+L)$. The action is free and properly discontinuous since the spatial translates move by integer multiples of $L$. Its differential is the identity, so both metric and time orientation descend. Write $b=\tau/L$. The pair $(T,[x])=(t-bx,x\bmod L)$ gives global product coordinates on the quotient: to invert, choose a representative $x$ and put $t=T+bx$; another choice changes the lift by a deck transformation. Thus the manifold is a cylinder for every $\tau$.

\begin{theorem}[Causal classification of the shifted cylinder]\label{f:shift-cylinder}
If $|\tau|>L$, every point lies on a closed timelike curve. If $|\tau|=L$, the cylinder is chronological but has closed null curves. If $|\tau|<L$, the function $T=t-(\tau/L)x$ is Cauchy temporal and the cylinder is globally hyperbolic.
\end{theorem}
\begin{proof}
The displacement $V=(\tau,L)$ has squared norm $L^2-\tau^2$. If it is timelike or null, straight segments of displacement $V$ (or $-V$ for future orientation) project to closed causal curves of that type. In the null case a future timelike loop would lift to a path with displacement $mV$ for an integer $m$. The strict inequality $|\dot x|<\dot t$ on timelike pieces gives $|\Delta x|<\Delta t$, impossible both for a null displacement and for $m=0$.

Suppose now $|b|<1$. The covector $\mathrm dT=\mathrm dt-b\,\mathrm dx$ has squared norm $b^2-1<0$ and is positive on future causal velocities. Along a future causal lift,
\[
 \dot T\geq(1-|b|)\dot t,\qquad |\dot x|\leq\dot t .
\]
Integrating between any two parameter values bounds the distance on the spatial circle by the increase of $T$ divided by $1-|b|$. Thus with $T$ as parameter the spatial path is Lipschitz into a complete circle. At a finite supremum of $T$ it would converge, giving an endpoint in the global coordinates $(T,[x])$. One could extend the path by increasing $T$ at fixed $[x]$, a timelike continuation. A finite infimum is excluded by the same past-endpoint argument. Therefore an inextendible path takes every value of $T$, exactly once by strict increase. Theorem~\ref{f:cauchy-splitting} applies to this smooth boundaryless quotient.
\end{proof}

\begin{figure}[htbp]\centering
\begin{tikzpicture}[x=1.13cm,y=1.02cm,>=Stealth,font=\small]
 \fill[lightpine] (0,0)--(-2.45,2.45)--(2.45,2.45)--cycle;
 \draw[pine] (-2.45,2.45)--(0,0)--(2.45,2.45);
 \draw[->,slate] (-2.6,0)--(3.5,0) node[right] {$x$};
 \draw[->,slate] (0,-.23)--(0,2.7) node[above] {$t$};
 \draw[dashed,slate] (1.4,0) node[below] {$L$}--(1.4,2.38);
 \draw[thick,->,brick] (0,0)--(1.4,.5) node[right=3pt] {spacelike};
 \draw[thick,->,ink] (0,0)--(1.4,1.4) node[right=3pt] {null};
 \draw[thick,->,pine] (0,0)--(1.4,2.18) node[above left=1pt] {timelike};
\end{tikzpicture}
\caption{For fixed $L$ the endpoint of the identification vector moves along the dashed line (here $\tau\geq0$). The light cone tests the straight segment. Ruling out other closed timelike paths at the null threshold requires the integrated cone inequality, not just the drawing.}\label{f:shift-cone-figure}
\end{figure}

The endpoint argument depends on which points and curves belong to the space. For instance in the open half-plane $x>0$, the future null ray $(t,-t)$ for $t<0$ stops at the omitted wall. The diamond from $(-2,1)$ to $(2,1)$ contains $(0,1/j)$, $j=1,2,\ldots$, but this sequence has no limit point there. In the closed half-plane the ray has an endpoint and can continue along the timelike wall if such curves are admitted. Results for boundaryless spacetimes do not automatically extend to regions with faces or corners; the permissible curves at the boundary and the meaning of inextendibility are part of a boundary version of the problem.

\begin{exercise}[Arrival after winding]
On the globally hyperbolic cylinder, determine the earliest $s>0$ for a causal signal from $[(0,0)]$ to $[(s,0)]$ that winds at least once. Can a timelike signal attain the same first arrival?

\emph{Solution.} A lift ending after $m\ne0$ windings has displacement $(s+m\tau,mL)$. The necessary inequality $s+m\tau\geq |m|L$ is sufficient by the straight segment. Minimizing $|m|L-m\tau$ over nonzero integers gives $L-|\tau|$. Equality is null; timelike paths require a strict inequality. The comparison measures a global travel time, not a different cone at the point of emission.
\end{exercise}

\begin{exercise}[Which spatial ends are harmless?]
Let $(\Sigma,h)$ be a connected Riemannian manifold without boundary. For the product metric $-\mathrm dt^2+h$ on $\mathbb R\times\Sigma$, show that $t$ is Cauchy temporal if and only if $\Sigma$ is complete for its Riemannian distance (the infimum of path lengths).

\emph{Hint.} Future causal paths satisfy $|\dot y|_h\leq\dot t$. In the complete case repeat the finite-endpoint argument of Theorem~\ref{f:shift-cylinder}. In the incomplete case choose a nonconvergent Cauchy sequence and a subsequence with successive distances less than $2^{-j-1}$. Join successive points by paths of lengths less than $2^{-j}$. Following them with unit spatial speed takes finite total time, but no spatial endpoint can exist. Extend into the past at fixed initial position; this inextendible causal path misses all sufficiently late slices.
\end{exercise}
